\subsection{Архимед. База}
%%%%%%%%%%%%%%%%%%%%%%%%%%%%%%%%%%%%%%%%%%%%%%%%%%%%
%%%%%%%%%%%%%%%%%%%%%%%%%%%%%%%%%%%%%%%%%%%%%%%%%%%%
 \z[{\begin{wrapfigure}[3]{r}{0.2\textwidth}%
		\vspace{-15pt}%
		\begin{flushright}%
 			\includegraphics[width=0.2\textwidth]{9.2.1}%
		\end{flushright}%
	\end{wrapfigure}%
	}]%
	{На границе раздела двух жидкостей плотности $\rho_1$ и $\rho_2$ плавает шайба плотности $\rho$ ($\rho_1 < \rho < \rho_2$). Высота шайбы $H$. Определите глубину $h$ её погружения во вторую жидкость.
	
	\vspace{2em}
	
	}
%
{$h = H \dfrac{\rho - \rho_1}{\rho_2 - \rho_1}$}
%%%%%%%%%%%%%%%%%%%%%%%%%%%%%%%%%%%%%%%%%%%%%%%%%%%%
\z[{\begin{wrapfigure}[4]{r}{0.2\textwidth}%
		\vspace{-15pt}%
		\begin{flushright}%
 			\includegraphics[width=0.2\textwidth]{9.2.2}%
		\end{flushright}%
	\end{wrapfigure}%
	}]%
	{С какой силой $F$ давит тяжелая палочка на дно водоема, если жестко связанный с палочкой пустотелый шарик радиуса $r$ погрузился в жидкость наполовину? Плотность жидкости $\rho$, длина палочки $l$.
		
	\vspace{2em}
	
	}
%
{$F = \dfrac{2}{3} \pi r^3 \rho g \left(1 + \dfrac{2r}{l}\right)$}
%%%%%%%%%%%%%%%%%%%%%%%%%%%%%%%%%%%%%%%%%%%%%%%%%%%% 
 \z[{\begin{wrapfigure}[4]{r}{0.15\textwidth}%
		\vspace{-35pt}%
		\begin{flushright}%
 			\includegraphics[width=0.17\textwidth]{9.2.3}%
		\end{flushright}%
	\end{wrapfigure}%
	}]%
	{Определите силу натяжения $T$ нити, связывающей два шарика объёма $10\,\text{см}^3$ каждый, если верхний шарик плавает, наполовину погрузившись в воду. Нижний шарик в три раза тяжелее верхнего.
		
	\vspace{2em}
	
	}
%
{$F = 1.2 \cdot 10^{-2}\,\text{Н}$}
%%%%%%%%%%%%%%%%%%%%%%%%%%%%%%%%%%%%%%%%%%%%%%%%%%%% 
 \z[{\begin{wrapfigure}[4]{r}{0.2\textwidth}%
		\vspace{-25pt}%
		\begin{flushright}%
 			\includegraphics[width=0.2\textwidth]{9.2.4}%
		\end{flushright}%
	\end{wrapfigure}%
	}]%
	{В цилиндрическом сосуде с водой плавает поплавок, к которому привязан груз массой $3\,\text{кг}$ и объёмом $1\,\text{л}$. Как и насколько изменится уровень воды $\Delta h$ в сосуде, если нить оборвётся и груз упадёт на дно? Площадь дна сосуда $50\,\text{см}^2$.
		
	\vspace{2em}
	
	}
%
{$\Delta h = -20\,\text{см}$}
%%%%%%%%%%%%%%%%%%%%%%%%%%%%%%%%%%%%%%%%%%%%%%%%%%%%	
 \z[{\begin{wrapfigure}[6]{r}{0.2\textwidth}%
		\vspace{+5pt}%
		\begin{flushright}%
 			\includegraphics[width=0.2\textwidth]{9.2.5}%
		\end{flushright}%
	\end{wrapfigure}%
	}]%
	{Два одинаковых бревна расположены в воде так, как показано на рисунке. Нижнее бревно привязано к вертикальной стенке тросами, составляющими с ней угол $45\degree$. Верхнее бревно наполовину погружено в воду. Определите плотность $\rho_{\text{др}}$ древесины.}
%
{$\rho = 2,\!3\,\text{г/см}^3$}